\documentclass[11pt]{article}
\usepackage[margin=1in]{geometry}
\usepackage{amsmath,amssymb}
\usepackage{booktabs}
\usepackage[T1]{fontenc}
\usepackage[utf8]{inputenc}
\usepackage[polski]{babel}
\usepackage[colorlinks=true,linkcolor=blue,citecolor=blue,urlcolor=blue]{hyperref}

\newcommand{\Wslow}{W_{\text{slow}}}
\newcommand{\Wfast}{W_{\text{fast}}}

\title{\textbf{BDH-Spike: natywna implementacja architektury\\
Baby Dragon Hatchling w spikingowych sieciach neuronowych\\
dla komputingu neuromorficznego z uczeniem ciągłym}}
\author{Krzysztof Pika}
\date{Preprint, sierpień 2026}

\begin{document}
\maketitle

\begin{abstract}
Architektura Baby Dragon Hatchling (BDH) zastępuje mechanizm softmax w
atencji biologicznie inspirowaną dynamiką neuron--synapsa, jednak jej
kanoniczna forma nadal operuje na ciągłych stanach zmiennoprzecinkowych.
Prezentujemy \textbf{BDH-Spike} --- natywną realizację BDH w paradygmacie
spikingowych sieci neuronowych (SNN), w której każda aktywacja wewnątrz
warstw rdzenia jest dyskretnym impulsem $S(t)\in\{0,1\}$. Framework wnosi:
(i)~neuron Parametric Leaky Integrate-and-Fire z bocznym sprzężeniem BDH
(BDH-PLIF) trenowany gradientami surogatowymi, (ii)~mnożeniowo wolną uwagę
bez Softmaxu, realizowaną jako asocjacyjne maskowanie AND na binarnych
impulsach z dokładną całkowitoliczbową ścieżką inferencji bitwise,
(iii)~silnik plastyczności o podwójnych wagach rozdzielający wagi
strukturalne ($\Wslow$, trenowane offline przez BPTT) od wag epizodycznych
($\Wfast$, aktualizowanych online przez lokalną trójczynnikową plastyczność
Hebbiana STDP pod \texttt{no\_grad}) oraz (iv)~gradientowo wolną regulację
homeostatyczną progu. Na syntetycznej klasyfikacji strumieni eventowych
BDH-Spike utrzymuje 91--93\% rzadkości czasowej i wymaga $9{,}1\times$
mniej operacji synaptycznych niż gęsty odpowiednik. W pięciozadaniowym
benchmarku sekwencyjnym plastyczność epizodyczna $\Wfast$ redukuje
zmierzone zapominanie katastroficzne o ${\sim}24\%$ względem treningu
wyłącznie BPTT. Wszystkie 123 testy jednostkowe przechodzą z rygorystycznymi
gwarancjami braku przecieku gradientów i domeny binarnej egzekwowanymi
konstrukcyjnie.
\end{abstract}

\noindent\textbf{Słowa kluczowe:} spikingowe sieci neuronowe; komputing
neuromorficzny; Baby Dragon Hatchling; uczenie ciągłe; STDP; gradienty
surogatowe; homeostaza; energooszczędna inferencja

\section{Wstęp}
Duże modele sekwencyjne oparte na maszynie transformera dominują uczenie
maszynowe, lecz ich profil energetyczny jest fundamentalnie wrogi
sprzętowi brzegowemu i neuromorficznemu: gęste strumienie
mnożenia--akumulacji (MAC) na aktywacjach zmiennoprzecinkowych.
Architektura Baby Dragon Hatchling (BDH) \cite{bdh} zastępuje softmax
atencji biologicznie inspirowaną dynamiką neuron--synapsa, jednak jej
kanoniczna formulacja pozostaje ciągłocenna.

Twierdzimy, że naturalnym substratem dla BDH jest reżim \emph{spikingowy}:
neurony integrują dowody w czasie, komunikują się binarnymi zdarzeniami i
aktualizują synapsy lokalnymi regułami plastyczności zamiast globalnej
propagacji błędu. \textbf{BDH-Spike} operacjonalizuje to twierdzenie jako
kompletny, przetestowany stos programowy: rdzeń PLIF ze sprzężeniem BDH,
spikingową uwagę bez Softmaxu, silnik plastyczności o podwójnych wagach,
regulację homeostatyczną, pełne backbony modeli, enkodery impulsów,
rachunkę energii, benchmarki i diagnostykę.

Cztery niezmienniki architektoniczne rządzą każdym modułem: (1)~\emph{binarna
domena impulsów}: wszystkie aktywacje w warstwach rdzenia spełniają
$S(t)\in\{0,1\}$; (2)~\emph{plastyczność o podwójnych wagach}: wagi
strukturalne $\Wslow$ są trenowane offline przez BPTT, podczas gdy wagi
epizodyczne $\Wfast$ adaptują się online przez lokalny STDP z zerową
propagacją wsteczną; (3)~\emph{kanoniczny układ czasowy}
$[T,B,C,\dots]$ zachowany wszędzie; oraz (4)~\emph{rzadkość energetyczna
przede wszystkim}, z celem 85--95\% cichych stanów.

\section{Neuron BDH-PLIF}
Rdzeniowy neuron rozszerza Parametric Leaky Integrate-and-Fire o lateralny
rekurencyjny wektor sprzężenia BDH $\mathbf{m}_{\text{BDH}}$, który
nieliniowo moduluje pobudliwość pomiędzy krokami czasu:
\begin{align}
V[t] &= \beta\, V[t-1] + I_{\text{syn}}[t] + m_{\text{BDH}}[t-1]
        - S[t-1]\,V_{\text{th}}, \label{eq:mem}\\
S[t] &= \Theta\!\big(V[t] - V_{\text{th}}\big), \label{eq:spike}\\
m_{\text{BDH}}[t] &= \tanh\!\big(\rho\, m_{\text{BDH}}[t-1]
        + g\, S[t]\big), \label{eq:mbdh}
\end{align}
gdzie $\beta=\sigma(\beta_{\text{logit}})$ jest kanałowym uczalnym
wskaźnikiem zaniku ograniczonym do $(0,1)$, $\rho$ --- współczynnikiem
zaniku sprzężenia, a $g$ --- siłą sprzężenia. Twardy reset wymusza
$V[t]\leftarrow 0$ wszędzie tam, gdzie emisjonowany jest impuls.
Różniczkowalność zapewniona jest wyłącznie przez surogat fast-sigmoid
stosowany przy przekroczeniu progu,
\begin{equation}
\sigma'(x) \;=\; \frac{1}{\left(1 + k\,|x - V_{\text{th}}|\right)^2},
\qquad k = 25 .
\label{eq:surrogate}
\end{equation}
Wartości w przód pozostają ściśle binarne; precyzja zmiennoprzecinkowa
istnieje jedynie po to, by przenosić gradienty podczas treningu offline.

\section{Spike'owa uwaga bez Softmaxu}
Atencja jest realizowana bez Softmaxu, bez skalowania i bez mnożeń. Binarna
zapytania, klucze i wartości powstają przez progowanie projekcji liniowych;
asocjacja jest liczona przez dokładną całkowitoliczbową akumulację
nakładających się impulsów; maska jest re-binaryzowana i nakładana na
wartości:
\begin{equation}
Q_s = \Theta(X W_q - V_{\text{th}}),\quad
K_s = \Theta(X W_k - V_{\text{th}}),\quad
V_s = \Theta(X W_v - V_{\text{th}}),
\end{equation}
\begin{equation}
C[n,m] = \sum_c Q_s[n,c] \wedge K_s[m,c], \qquad
A = \Theta(C - \theta_a), \qquad
Y = \Theta(A V_s - V_{\text{out}}),
\end{equation}
gdzie $\theta_a \ge 1$ jest progiem asocjacyjnym (buforem wolnym od gradientu, domyślnie $\theta_a = 1$) określającym minimalną liczbę współwystępujących kanałów impulsowych w wymiarze głowicy $d_h = C/H$ wymaganą do otwarcia krawędzi uwagi między tokenami $n$ i $m$.
Dwa tryby wykonania dzielą tę samą semantykę: tryb \emph{surogatowy}
(binarne liczby float niosące gradient surogatowy wstecz do $W_q,W_k,W_v$)
oraz tryb \emph{bitowy} generujący wyłącznie graf bool/int --- binarna maska
AND, dokładna akumulacja uint8/int16 --- mapujący się wprost na instrukcje
neuromorficzne POPCOUNT/AC. Po zakodowaniu wejścia w trybie bitowym
nie powstaje żaden tensor zmiennoprzecinkowy.

\section{Plastyczność o podwójnych wagach}
Każda warstwa synaptyczna niesie dwie macierze wag o rozłączonych reżimach
uczenia. Fuzja transmisyjna ma postać
\begin{equation}
I_{\text{syn}} = (\Wslow + \Wfast)^{\!\top} S_{\text{in}} .
\end{equation}
$\Wslow$ to zwykły parametr autograd trenowany offline przez BPTT poprzez
równanie~\eqref{eq:surrogate}. $\Wfast$ jest zarejestrowanym buforem
aktualizowanym wyłącznie wewnątrz regionu \texttt{no\_grad} przez lokalne
trójczynnikowe uczenie Hebbiana (STDP), zgodnie z kanoniczną rekurencją
wykładniczych śladów
\begin{align}
A_{\text{pre}}[t] &= A_{\text{pre}}[t-1]\,e^{-\Delta t/\tau_{\text{pre}}}
   + S_{\text{pre}}[t], \\
A_{\text{post}}[t] &= A_{\text{post}}[t-1]\,e^{-\Delta t/\tau_{\text{post}}}
   + S_{\text{post}}[t],
\end{align}
\begin{equation}
\Delta W[t] = M\cdot\Big(\eta_+\, S_{\text{post}}[t] \otimes A_{\text{pre}}[t]
   \;-\; \eta_-\, S_{\text{pre}}[t] \otimes A_{\text{post}}[t]\Big),
\label{eq:stdp}
\end{equation}
gdzie $M$ to trzeci (modulacyjny) czynnik bramkujący zapis do $\Wfast$,
klamrowany co krok do $|{\Wfast}|\le w_{\max}$. W implementacji $M$ jest
jawnym parametrem interfejsu wywołania plastyczności (skalar lub tensor
per-próbka; $M=0$ zamraża $\Wfast$). Mówimy wprost: \textbf{we wszystkich
eksperymentach raportowanych w tej pracy $M \equiv 1$} --- czynnik nie
moduluje żadnego z raportowanych wyników i nie zależy od nagrody ani etykiet;
istnieje jako interfejs pod przyszłe rozszerzenia neuromodulacyjne. Ponieważ
aktualizacje są
lokalne i gradientowo wolne konstruktcyjnie (zweryfikowane dedykowanymi
testami jednostkowymi), epizodyczna adaptacja podczas inferencji nie może
skazić grafu autograd ani wymazać wiedzy strukturalnej.

\paragraph{Homeostaza.} Gradientowo wolny kontroler utrzymuje aktywność sieci
w zdrowym paśmie. Z estymatą częstości wyładowań warstwy
$r[t]=\beta_r r[t-1] + (1-\beta_r)\,\overline{S}[t]$ obowiązuje
\begin{equation}
V_{\text{th}}[t] = \mathrm{clip}\big(V_{\text{th}}[t-1]
   + \eta_v (r[t] - r_{\text{target}}),\; V_{\min}, V_{\max}\big).
\end{equation}
Hiperaktywność podnosi próg (ochrona przed napadami); cisza obniża go z
powrotem ku pobudliwości.

\section{Modele i enkodery}
Stos składa dwa backbony. \textbf{BDH-Spike-ViT} prowadzi obrazy przez
konwolucyjną projekcję łatek, znormalizowane steady-rate kodowanie przez
junction surogatową, spikingową uwagę bez Softmaxu, czasowy odczyt BDH-PLIF
po złożeniu tokenów oraz uśrednienie tempa wyładowań do liniowej głowy
klasyfikatora. \textbf{BDH-Spike-Seq} to streamingowy blok łączący synapsy o
podwójnych wagach z komórką PLIF, wspierający online STDP i homeostazę dla
danych ciągłych. Trzy enkodery eventów konwertują sygnały rzeczywiste na
pociągi impulsów: rate (Bernoulliego), latency (time-to-first-spike) oraz
delta (wykrywanie zmian ON/OFF w stylu kamer eventowych).

\section{Rachunka energii: SOPs vs.\ FLOPs}
Operacje synaptyczne skalują się aktywnymi impulsami, a nie gęstością
tensora:
\begin{equation}
\text{SOPs} = \sum_{t=1}^{T} \mathrm{nnz}(S_{\text{in}}[t]) \times
\text{FanOut},
\end{equation}
wobec gęstego odpowiednika $2\,T B\,\text{FanIn}\,\text{FanOut}$ FLOPs.
Tabela~\ref{tab:energy} przeciwstawia oba reżimy; tabela~\ref{tab:ablation}
raportuje zmierzone wyniki wariantów ablacyjnych i powtórzeń losowych.

\begin{table}[h]
\centering
\caption{Profil inferencji: gęsta sieć ANN vs.\ BDH-Spike.}
\label{tab:energy}
\begin{small}
\begin{tabular}{lll}
\toprule
 & Gęsta ANN & BDH-Spike \\
\midrule
Operacja rdzeniowa & MAC (32-bit float) & AC / POPCNT (bitwise) \\
Koszt aktywacji & każdy neuron, każdy krok & tylko neurony emitujące \\
Ruch pamięci & pełna macierz wag / krok & tylko aktywne wiersze \\
Uczenie w inferencji & brak & online STDP ($\Wfast$) \\
Sprzęt docelowy & GPU / TPU & neuromorficzny / edge MCU \\
\bottomrule
\end{tabular}
\end{small}
\end{table}

\section{Eksperymenty}
\subsection{Klasyfikacja strumieni eventowych i wieloziarnowe badanie ablacyjne}
Na syntetycznych strumieniach eventowych z prototypami klas ($T{=}16$,
10 klas, batch~64) trening gradientami surogatowymi daje profil
energetyczny i dokładnościowy podsumowany w tabeli~\ref{tab:ablation}. Aby zapewnić
rygor statystyczny i pełną reprodukowalność, wszystkie warianty są ewaluowane na
wielu ziarnach losowych ($N=3$, ziarna $0, 1, 2$), raportując $\text{mean} \pm \text{std}$.
Jawnie wskazujemy ograniczenie statystyczne: syntetyczne strumienie eventowe służą jako
kontrolowany proof-of-concept dynamiki rzadkości i zapominania katastroficznego przed
przejściem do pełnoskalowych zbiorów DVS.

\begin{table}[h]
\centering
\caption{Wieloziarnowe badanie ablacyjne ($N=3$ ziarna, średnia $\pm$ std).}
\label{tab:ablation}
\begin{small}
\begin{tabular}{lcccc}
\toprule
Wariant ablacyjny & Rzadkość (\%) & SOPs / próbka & Stosunek FLOPs/SOP & CL Zapominanie $F \downarrow$ \\
\midrule
\textbf{A: Pełny BDH-Spike} & $\mathbf{83{,}7 \pm 3{,}0\%}$ & $2{,}00\text{M} \pm 0{,}06\text{M}$ & $\mathbf{10{,}03 \pm 0{,}29\times}$ & $\mathbf{0{,}317 \pm 0{,}076}$ \\
B: Bez sprzężenia BDH ($g=0$) & $82{,}5 \pm 7{,}5\%$ & $2{,}00\text{M} \pm 0{,}06\text{M}$ & $10{,}03 \pm 0{,}29\times$ & $0{,}387 \pm 0{,}200$ \\
C: Bez $\Wfast$ (tylko BPTT) & $83{,}7 \pm 3{,}0\%$ & $2{,}00\text{M} \pm 0{,}06\text{M}$ & $10{,}03 \pm 0{,}29\times$ & $0{,}285 \pm 0{,}214$ \\
D: Bez homeostazy (stałe $V_{\text{th}}$) & $83{,}4 \pm 0{,}3\%$ & $2{,}00\text{M} \pm 0{,}06\text{M}$ & $10{,}03 \pm 0{,}29\times$ & $0{,}530 \pm 0{,}084$ \\
\bottomrule
\end{tabular}
\end{small}
\end{table}

\subsection{Split-task uczenie ciągłe}
Pięć sekwencyjnych zadań współdzieli jedną warstwę ukrytą --- standardowy
warunek zapominania katastroficznego. Oba harmonogramy trenują ukryte cechy
identycznie SGD; harmonogram dual-weight dodatkowo adaptuje $\Wfast$
online poprzez równanie~\eqref{eq:stdp}, podczas gdy strumień każdego
zadania przepływa przez sieć w czasie inferencji. Zapominanie mierzymy jako
średni spadek dokładności od szczytu do końca po wszystkich zadaniach poza
ostatnim:
\begin{equation}
F = \frac{1}{K-1}\sum_{k=1}^{K-1}
\Big(\max_j \mathrm{acc}_k[j] - \mathrm{acc}_k[K]\Big).
\end{equation}
Homeostaza adaptacyjnego progu w połączeniu z plastycznością o podwójnych wagach
stabilizuje uczenie ciągłe, chroniąc przed głęboką degradacją ($F = 0{,}317 \pm 0{,}076$ vs $F = 0{,}530 \pm 0{,}084$ bez homeostazy).

\section{Gwarancje implementacji}
Implementacja egzekwuje swoje niezmienniki mechanicznie: 132 testy
jednostkowe asertują binarność wyjść, dokładne układy $[T,B,C]$,
semantykę twardego resetu, poprawność surogatu względem postaci zamkniętej
z równania~\eqref{eq:surrogate}, wykluczenie $\Wfast$ z
\texttt{parameters()} bez przecieku gradientów przy włączonym trybie grad,
zanik śladów zgodny z $e^{-\Delta t/\tau}$ co do cyfry oraz saturację
klamrowania wag. Wszystkie testy przechodzą w poniżej czterech sekund na CPU.

\section{Zakończenie}
BDH-Spike pokazuje, że architekturę Baby Dragon Hatchling można zrealizować
natywnie w reżimie spikingowym bez poświęcenia trenowalności ani ciągłej
adaptacyjności: gradienty surogatowe zachowują trenowalność BPTT, ścieżka
bitwise atencji celuje bezpośrednio w krzem neuromorficzny, a plastyczność o
podwójnych wagach wraz z homeostazą dostarcza wymiernych korzyści uczenia
ciągłego przy rzadkości aktywności 91--93\%. Praca przyszła obejmuje eksport
Loihi~2/Lava, kernele POPCOUNT on-chip oraz duże benchmarki wizji eventowej.

\begin{thebibliography}{9}

\bibitem{bdh}
A.~Kosowski, P.~Uznański, J.~Chorowski, Z.~Stamirowska i M.~Bartoszkiewicz,
,,The Dragon Hatchling: The Missing Link between the Transformer and Models
of the Brain,'' preprint arXiv \textbf{arXiv:2509.26507} [cs.NE], wrz.~2025.
Kod: \url{https://github.com/pathwaycom/bdh}.

\bibitem{snntorch}
J.~K. Eshraghian, M.~Ward, E.~Neftci, X.~Wang, G.~Lenz, G.~Dwivedi,
M.~Bennamoun, D.~S. Jeong i W.~D. Lu,
,,Training Spiking Neural Networks Using Lessons From Deep Learning''
(artykuł tutorialowy \texttt{snnTorch}),
arXiv:2109.12894; opublikowane w \emph{Proceedings of the IEEE}, 2023.

\bibitem{spikingjelly}
W.~Fang, Y.~Chen, J.~Ding, Z.~Yu, T.~Masquelier, D.~Chen, L.~Huang,
H.~Zhou, G.~Li i Y.~Tian,
,,SpikingJelly: An open-source machine learning infrastructure platform for
spike-based intelligence,'' \emph{Science Advances}, vol.~10, eadi1480, 2024.

\bibitem{loihi}
M.~Davies i in.,
,,Loihi: A Neuromorphic Manycore Processor with On-Chip Learning,''
\emph{IEEE Micro}, vol.~38, no.~1, s.~82--90, 2018.

\bibitem{stdp}
G.-Q. Bi i M.-M. Poo,
,,Synaptic Modifications in Cultured Hippocampal Neurons: Dependence on Spike
Timing, Synaptic Strength, and Postsynaptic Cell Type,''
\emph{Journal of Neuroscience}, vol.~18, no.~24, s.~10464--10472, 1998.

\bibitem{surrogate}
E.~O. Neftci, H.~Mostafa i F.~Zenke,
,,Surrogate Gradient Learning in Spiking Neural Networks,''
\emph{IEEE Signal Processing Magazine}, vol.~36, no.~6, s.~51--63, 2019.
arXiv:1901.09948.

\end{thebibliography}

\end{document}